% =============================================================================
% cap2_hardware/05_tipologias.tex
% Seção 2.5 -- Tipologias de câmera (v2.2 — segunda compactação)
% =============================================================================
\section{Tipologias de câmera}
\label{sec:tipologias-arquiteturais}

Consolidados o conjunto de requisitos (Seção~\ref{sec:requisitos-tipologias}) e as condições de contorno físicas (Seção~\ref{sec:acondicionamento-geometria}) e energéticas (Seção~\ref{sec:energia-conectividade}), torna-se possível percorrer comparativamente as quatro tipologias arquiteturais hoje disponíveis para aquisição autônoma em campo. A discussão é organizada como \emph{escada tecnológica} ordenada por compatibilidade crescente com os requisitos desejáveis. A leitura não é hierárquica no sentido econômico: cada degrau resolve parte do problema e introduz novos \emph{trade-offs}. A escolha consolidada para o Campus~2, articulada na Seção~\ref{sec:decisao-arquitetural}, fixa um único degrau como tipologia operacional e reserva os demais a finalidades experimentais e de evolução futura.

\paragraph*{Câmeras IP comerciais domésticas.} O primeiro degrau é o de maior acessibilidade: o reaproveitamento de \textbf{câmeras IP de consumo} (IPCAMs), originalmente projetadas para vigilância residencial com energia elétrica contínua e Wi-Fi estável, adaptadas para operação autônoma. Essa adaptação envolve três intervenções complementares. A primeira é a substituição da alimentação de rede por bancos de bateria recarregável (\emph{powerbanks}) dimensionados para múltiplos dias de operação, com gravação local em cartão SD industrial em substituição ao armazenamento em nuvem do fabricante. A segunda é a modificação do \emph{firmware}: muitas IPCAMs são distribuições Linux embarcadas (kernel customizado sobre \emph{BusyBox}) cujo acesso, por interfaces de \emph{debug} ou \emph{firmware} alternativo, permite desabilitar rádios não utilizados (Wi-Fi, Bluetooth), suprimir serviços de telemetria e ajustar a política de gerenciamento de energia ao caso de uso real. Iniciativas comunitárias de \emph{firmware} aberto --- com destaque para o projeto \textbf{OpenIPC}, voltado a câmeras IP com SoCs HiSilicon, Xiongmai, SigmaStar, GoKe e Ingenic --- têm tornado tecnicamente acessível a troca do \emph{firmware} de fábrica por imagens enxutas, auditáveis e configuráveis~\cite{openipc2025}, ao custo de compatibilidade verificada modelo a modelo. A terceira intervenção é a integração mecânica ao abrigo, conforme a Seção~\ref{sec:acondicionamento-geometria}. A objeção crítica a esse degrau está ligada ao primeiro requisito mínimo: quando a detecção de movimento é implementada por \emph{software}, o subsistema de imagem precisa permanecer permanentemente ativo, drenando rapidamente a bateria mesmo sem captura efetiva. Modelos com sensor PIR de \emph{hardware} nativo amenizam o problema, mas o consumo em modo ocioso permanece superior ao de arquiteturas projetadas para operação autônoma desde a origem~\cite{velasco2024smartcameratraps}. Para o pipeline do Capítulo~\ref{cap:pipeline}, esse degrau implica receber quadros vindos de um codec proprietário, com pouca margem para metadados de inferência embarcada e taxa de \emph{ghost frames} potencialmente elevada quando o disparo é puramente por software.

\paragraph*{Câmeras de trilha (\emph{trail cameras}).} O segundo degrau resolve o ponto fraco anterior por desenho. As \emph{trail cameras}, herdeiras de três décadas de tradição em ecologia de fauna~\cite{bruce2025largescale,cordier2022africa,swanson2015snapshot}, são equipamentos integrados que reúnem nativamente sensor PIR de hardware, vedação ambiental classe IP66 ou IP67, iluminação infravermelha próxima embarcada e autonomia superior a meses de operação com baterias alcalinas ou de lítio. Esses atributos as posicionam, em robustez mecânica e autonomia energética, como referência do segmento. O custo dessa maturidade é a \textbf{arquitetura fechada}: o \emph{firmware} é proprietário e não admite injeção de algoritmos customizados na borda; o pipeline de imagens é encerrado no ciclo gravação $\to$ cartão SD $\to$ coleta manual; e mesmo modelos com transmissão 4G nativa operam sob protocolos e nuvens proprietárias, dificultando a integração ao restante da arquitetura. A \emph{trail camera} cumpre todos os requisitos mínimos mas falha sistematicamente nos desejáveis: não há \emph{Edge AI} possível e a sincronização, quando existe, é cativa. Consequentemente, nenhuma etapa do pipeline do Capítulo~\ref{cap:pipeline} pode ser deslocada para a borda, e todas as decisões de detecção, classificação e re-ID precisam ocorrer a posteriori em servidor.

\paragraph*{Microcontroladores com câmera embarcada (ESP32-CAM e similares).} O terceiro degrau é o do \emph{hardware} aberto. Plataformas como o ESP32-CAM reúnem custo unitário irrisório, consumo extremamente baixo em modo ocioso, suporte nativo a PIR de \emph{hardware} por GPIO, comunicação Wi-Fi e gravação em cartão SD em um único módulo. São tipicamente utilizadas em prototipagem, em redes densas de baixo custo e em pontos onde a perda do equipamento é aceitável~\cite{velasco2024smartcameratraps}. O \emph{trade-off} é duplo: a qualidade ótica do sensor embarcado é marginal para tarefas exigentes como a reidentificação individual (iluminação, foco e faixa dinâmica insuficientes para os padrões finos de pelagem) e a capacidade computacional do microcontrolador, ainda que admita modelos diminutos via TinyML, é insuficiente para hospedar detectores generalistas como o \modeloMD. O módulo opera, portanto, essencialmente como conduto de imagens, adequado a tarefas binárias de presença e a prototipagem, mas inadequado como solução final em um sistema cuja finalidade é a reidentificação individual em ambiente urbano. Para o pipeline, esse degrau não modifica a divisão borda-servidor frente à \emph{trail camera} e adiciona limitação de qualidade óptica de entrada que se tornaria fator restritivo da própria etapa de re-ID.

\paragraph*{Computadores em placa única com aceleradores neurais (SBC + NPU).} O quarto degrau é o paradigma mais alinhado aos requisitos desejáveis. Computadores em placa única (\emph{Single Board Computers}, SBC), como o Raspberry Pi, são empregados há mais de uma década em monitoramento ecológico de baixo custo, com casos documentados em ninhos de aves, comportamentos crípticos de pequenos vertebrados e ambientes hostis~\cite{sargent2021raspberrypi}. Sem aceleração neural dedicada, contudo, esses sistemas operam no limite do orçamento computacional, com latência e consumo incompatíveis com inferência contínua~\cite{tiddeman2023lowpowerhardware}. A combinação do SBC com uma \textbf{Unidade de Processamento Neural} (NPU) --- a exemplo do Google Coral Edge TPU ou do Hailo-8 --- altera essa equação: torna-se possível executar localmente detectores generalistas como o \modeloMD~e classificadores multi-espécie como o \modeloSN, transmitindo apenas metadados estruturados ou quadros previamente filtrados~\cite{pishdast2025smartcameratraps,velasco2024smartcameratraps,mulero2025addressing}. O preço é um custo unitário significativamente superior e dimensionamento energético mais robusto, com painéis solares e bancos de bateria adequadamente especificados conforme a Seção~\ref{sec:energia-conectividade}. Para o pipeline, este é o único degrau que efetivamente habilita uma divisão de trabalho borda-servidor: detecção e classificação multi-espécie podem ocorrer no campo, descartando o ruído biológico antes da transmissão, e somente quadros pré-filtrados chegam ao servidor para reidentificação individual.

A leitura conjunta dos quatro degraus prepara a Seção~\ref{sec:decisao-arquitetural}: nenhuma das tipologias satisfaz simultaneamente todos os requisitos mínimos e desejáveis sob o orçamento típico de um projeto de graduação operado em ambiente cooperativo, e a escolha operacional resulta de compromisso entre custo, disponibilidade no comércio nacional, simplicidade de manejo e controle do nó de campo. As câmeras IP comerciais adaptadas, em particular sob \emph{firmware} aberto, oferecem o melhor desse compromisso para os dez pontos do Campus~2; as \emph{trail cameras} permanecem como referência para validação experimental em condições severas de energia; o SBC com NPU mantém-se como horizonte aspiracional para evoluções futuras com processamento de borda; e os microcontroladores ficam restritos a prototipagem ou a redes densas com finalidade exploratória. A Seção~\ref{sec:decisao-arquitetural} articula essa síntese em uma decisão arquitetural concreta para o Campus~2.
